\begin{mainquestion}{training}{Required EE101A Training}
\section*{Soldering Training}
Always wear safety glasses when soldering. Melted solder can ``pop,'' ejecting a small droplet of molten metal. If the drop of liquid metal falls on your skin, you may feel a slight sting. If the drop of liquid metal falls into your eye, you will have a very bad day.

When powering up a new soldering iron for the first time, promptly coat the tip with solder to prevent oxidation. The soldering tip usually has a copper core with an iron surface. If this surface comes in contact with oxygen (air) at elevated temperatures, it will oxidize to form iron oxide (``rust''). Iron oxides are thermal insulators that impede heat flow from the soldering tip.

Forming a protective layer of solder on the tip surface is called \emph{tinning the tip}. Please note that the soldering iron may take a couple of minutes to reach its operating temperature.

Please carefully read all of the information at the link below:

\begin{center}
\url{https://sites.google.com/stanford.edu/soldering-internal/learning?authuser=0}
\end{center}

\begin{partquestion}{solder}
{Please submit a screenshot of the completion screen at the end of the soldering quiz. The quiz can be found at the link above.}

\begin{answer}
\end{answer}
\end{partquestion}
\section*{Safety Training}
In order to gain access to the lab, you will need to complete the \href{https://drive.google.com/file/d/1UYZG5UO8G1NxXx5_a6gG9zF61goSwH3i/view}{EE Basic Lab Safety Training} and pass the \href{https://docs.google.com/forms/d/e/1FAIpQLSeOOwr2M8SxW7_CySm2XszzfyHqjzxqcVB9zLIyDYSYWiBs8A/viewform}{Lab Safety Quiz}.

\begin{partquestion}{safety}
{Please submit a screenshot of the completion screen at the end of the lab safety quiz.}

\begin{answer}
\end{answer}
\end{partquestion}


\end{mainquestion}